\section{本章小结}
\label{sec:fvec-summary}

本章开头的问题是：把函数看作向量$[\ldots,f(x),\ldots]$时，有限维线性代数的语言保留了什么，又在哪里失效。表~\ref{tab:fvec-summary}按本章的顺序汇总了答案。

\begin{BookInlineTable}
  \centering
  \begin{tabular}{@{}p{18mm}p{36mm}p{48mm}@{}}
    \toprule
    概念 & 有限维向量 & 无限维函数空间\\
    \midrule
    线性运算 & 逐分量计算 & 逐点计算，规则相同\\
    维数 & 分量个数 & 与输入点数无关，要看函数集合\\
    取值与坐标 & 标准基下坐标就是分量 & 取值一般不是有限组合的坐标\\
    基 & 有限个，有限组合表示全部 & 代数基只允许有限组合，通常写不出\\
    无穷和 & 用不到 & 收敛、范围、唯一性都依赖误差\\
    逐分量与统一误差 & 两种收敛等价 & 两种收敛不同\\
    有界序列 & 总有收敛子列 & 可以没有收敛子列\\
    \bottomrule
  \end{tabular}
  \caption{有限维向量的语言在无限维函数空间中的保留与失效。}
  \label{tab:fvec-summary}
\end{BookInlineTable}

线性运算的部分完整地保留下来。逐点加法与数乘使全体函数$X\to\mathbb K$成为线性空间，数列与有限维向量都是它的特例。连续函数、限次多项式、收敛到零的数列和有限支撑数列都通过了子空间核验。函数作为元素、它在某点的值、它相对于某组基的坐标是三样不同的东西；乘法和求导也不是每个函数空间自带的运算。

维数仍按张成整个空间需要多少方向来计算，与定义域的点数无关。有限输入集合上的函数以指示函数为基，坐标就是函数值。$P_m$定义在整条实轴上，却只有$m+1$维，取值列与系数列由一个可逆矩阵互相决定。全体多项式、全体数列含有任意多个线性无关的元素，因而是无限维的。其中每个元素可能只用到有限个方向，但没有一组固定的有限方向能表示所有元素。

无限维带来两个困难。其一，代数基只允许有限组合，要表示常数列或$e^x$这样的元素就须借助级数，而级数是否收敛、能表示哪些元素、系数是否唯一，都取决于选用的误差。其二，任何无限维赋范空间中，Riesz引理都给出两两距离至少$\tfrac12$的单位向量列，所以正半径闭球都不紧，这一结论不需要空间完备。

这两个困难分别指向后面两章。第~\ref{chap:function-normed-spaces}章用范数规定整体误差，比较不同范数下的收敛，并用完备性保证柯西列的极限留在空间内；第~\ref{chap:function-inner-product-spaces}章再加入内积，用正交投影和正交展开给出收敛的无限表示。闭球不紧则提醒我们：在无限维中求极值或抽收敛子列，需要额外的紧性条件或更弱的收敛概念。
